\chapter{Opioid medicines (how they work)}

\section*{Definition and Overview}
Opioid medicines are a class of powerful analgesic agents derived from the opium poppy (\textit{Papaver somniferum}) or synthesised chemically to mimic natural opiates. They are primarily prescribed for the management of moderate to severe acute pain, such as post-operative pain or pain from major trauma, and for some types of severe chronic pain, including cancer-related pain. Opioids exert their pharmacological effects by binding to specific receptors in the central and peripheral nervous systems, thereby altering the perception of and emotional response to pain. While highly effective, their use carries significant risks of tolerance, physical dependence, and addiction, necessitating careful clinical management and monitoring.

\section*{Epidemiology}
The use and prescribing patterns of opioid medicines have shifted significantly over the past few decades. In Australia and other Western countries, opioid prescriptions rose dramatically in the late 1990s and 2000s, leading to a public health crisis characterised by high rates of prescription opioid misuse, dependency, and accidental overdose. Recent regulatory changes in Australia, such as the introduction of real-time prescription monitoring (e.g., SafeScript) and the rescheduling of codeine to prescription-only, have aimed to reduce inappropriate prescribing. Despite these measures, pharmaceutical opioids remain a leading cause of drug-induced deaths in Australia, with middle-aged adults (specifically those aged 40 to 60) and individuals in regional or lower socioeconomic areas experiencing disproportionately higher rates of opioid-related harm.

\section*{Aetiology and Pathophysiology}
Opioids work by interacting with the body's endogenous opioid system, which regulates pain, reward, and addictive behaviours.
\begin{itemize}
    \item \textbf{Receptor binding:} Opioids act as agonists, binding to specific G-protein coupled receptors—primarily Mu ($\mu$), but also Kappa ($\kappa$) and Delta ($\delta$) opioid receptors—located on neuronal membranes in the brain, spinal cord, and gastrointestinal tract.
    \item \textbf{Inhibition of neurotransmitter release:} Binding of opioids to presynaptic receptors inhibits the entry of calcium ions, preventing the release of excitatory neurotransmitters (such as glutamate and substance P) that transmit pain signals.
    \item \textbf{Hyperpolarisation of postsynaptic neurons:} Opioids activate postsynaptic potassium channels, causing potassium ions to exit the cell. This hyperpolarises the postsynaptic membrane, making it less likely to generate an action potential in response to pain stimuli.
    \item \textbf{Descending pathway modulation:} In the brainstem, opioids inhibit GABAergic interneurons, thereby disinhibiting descending pathways that project to the spinal cord to suppress pain transmission at the dorsal horn.
    \item \textbf{Dopaminergic reward pathway stimulation:} By inhibiting GABA release in the ventral tegmental area, opioids increase dopamine release in the nucleus accumbens, producing feelings of euphoria and reinforcing the drive to take the drug, which underlies addiction.
\end{itemize}

\section*{Clinical Features}
The clinical effects of opioids include both therapeutic actions (analgesia) and a range of characteristic adverse effects across multiple organ systems.
\begin{itemize}
    \item \textbf{Analgesia:} Profound relief of moderate to severe pain, reducing both the physical sensation and the emotional distress associated with pain.
    \item \textbf{Central nervous system depression:} Drowsiness, sedation, cognitive impairment, and a feeling of calm or euphoria.
    \item \textbf{Respiratory depression:} Decreased sensitivity of the brainstem to carbon dioxide, leading to a reduced respiratory rate and depth, which is the primary cause of death in overdose.
    \item \textbf{Gastrointestinal hypomotility:} Severe constipation (opioid-induced constipation), nausea, vomiting, and delayed gastric emptying due to binding to enteric receptors.
    \item \textbf{Pupillary miosis:} Constriction of the pupils (pinpoint pupils) due to stimulation of the parasympathetic pathway of the oculomotor nerve.
    \item \textbf{Pruritus:} Itching, particularly around the nose and face, caused by systemic histamine release (especially common with morphine).
    \item \textbf{Endocrine effects:} Long-term use can suppress the hypothalamic-pituitary-gonadal axis, leading to reduced testosterone or oestrogen, resulting in sexual dysfunction, fatigue, and osteoporosis.
\end{itemize}

\section*{Differential Diagnosis}
While opioids are unique in their mechanism, the clinical presentation of opioid toxicity or their pharmacological profile must be distinguished from other states of central nervous system (CNS) depression or pain management strategies.
\begin{itemize}
    \item \textbf{Other sedative-hypnotic toxicities:} Poisoning by benzodiazepines, barbiturates, or alcohol, which cause CNS and respiratory depression but typically present with normal or dilated pupils (unless complicated by severe hypoxia).
    \item \textbf{Non-opioid analgesic profiles:} Non-steroidal anti-inflammatory drugs (NSAIDs) and paracetamol, which act through peripheral inflammatory pathways or central cyclooxygenase inhibition, lacking CNS and respiratory depressive effects.
    \item \textbf{Neuropathic pain agents:} Gabapentinoids (such as pregabalin and gabapentin) and tricyclic antidepressants, which modulate calcium channels or monoamine reuptake and do not bind to opioid receptors.
    \item \textbf{Endocrine or metabolic encephalopathy:} Hypoglycaemia, myxoedema coma, or hypercapnia, which present with altered consciousness but are distinguishable by blood chemistry and clinical history.
\end{itemize}

\section*{When to Seek Medical Care}
Patients taking prescribed opioids should consult their doctor if they experience inadequate pain control, severe constipation, significant drowsiness, or signs of physical dependence. Immediate medical attention is vital if signs of overdose occur.

\textbf{Call 000 (or local emergency services) for:}
\begin{itemize}
    \item Slow, shallow, or irregular breathing (fewer than 8 to 10 breaths per minute)
    \item Extreme somnolence, difficulty waking up, or unresponsiveness
    \item Pinpoint pupils in an unresponsive individual
    \item Blue or pale lips, face, or fingernails (cyanosis)
    \item Choking, snorting, or gurgling sounds (indicating partial airway obstruction or pulmonary oedema)
\end{itemize}

\section*{Diagnosis and Investigations}
Assessing the action, concentration, and safety of opioid medicines involves clinical examination, therapeutic monitoring, and toxicology.
\begin{itemize}
    \item \textbf{Clinical assessment:} Evaluation of pain scores, respiratory rate, level of consciousness, and pupillary size to monitor therapeutic efficacy and toxicity.
    \item \textbf{Urine drug screening:} Qualitatively detects the presence of opioids and their metabolites to monitor adherence or screen for illicit drug use.
    \item \textbf{Serum toxicology assays:} Quantitative measurement of specific opioids (e.g., paracetamol/codeine combinations) in emergency settings to guide management.
    \item \textbf{Arterial blood gas (ABG) analysis:} Assesses carbon dioxide retention (hypercapnia) and respiratory acidosis in cases of suspected respiratory depression.
    \item \textbf{Electrocardiogram (ECG):} Recommended for patients on high-dose methadone to monitor for QT interval prolongation, which increases the risk of Torsades de Pointes.
\end{itemize}

\section*{Management}
Opioid management requires careful titration, minimisation of adverse effects, and structured plans for cessation.
\begin{itemize}
    \item \textbf{Structured prescribing:} Restricting opioids to the lowest effective dose for the shortest possible duration, with clear treatment goals and a pre-defined exit strategy.
    \item \textbf{Opioid rotation:} Switching to a different opioid drug or administration route if the current agent causes intolerable side effects or inadequate analgesia.
    \item \textbf{Laxative therapy:} Proactive prescribing of stimulant laxatives (such as senna) and stool softeners (such as docusate) to manage opioid-induced constipation, as tolerance to constipation does not develop.
    \item \textbf{Antidote administration (Naloxone):} Emergency administration of naloxone, a competitive opioid receptor antagonist, to reverse life-threatening respiratory depression in overdose.
    \item \textbf{Adjuvant medications:} Combining opioids with paracetamol, NSAIDs, or neuropathic agents to allow opioid dose reduction (multimodal analgesia).
    \item \textbf{Opioid taper plans:} Gradual reduction of the opioid dose under medical supervision to prevent severe withdrawal symptoms in dependent patients.
\end{itemize}

\section*{Complications}
\begin{itemize}
    \item Severe, life-threatening respiratory depression leading to hypoxic brain injury or death
    \item Opioid-induced hyperalgesia, a paradoxical state where chronic opioid use increases sensitivity to pain
    \item Physical dependence and severe withdrawal syndrome upon abrupt cessation or dose reduction
    \item Opioid use disorder (addiction), characterised by compulsive use despite physical, social, or psychological harm
    \item Chronic endocrine dysfunction, including hypogonadism, infertility, erectile dysfunction, and adrenal suppression
    \item Increased risk of falls and fractures, particularly in elderly patients, due to sedation and orthostatic hypotension
\end{itemize}

\section*{Prognosis and Outcomes}
The short-term prognosis for pain relief with opioid medicines is generally excellent when used for acute, post-operative, or terminal cancer pain. However, the long-term prognosis for chronic non-cancer pain managed with opioids is less favourable; evidence indicates that long-term opioid therapy rarely leads to sustained functional improvement and is frequently associated with diminished quality of life, opioid-induced hyperalgesia, and dependency. For individuals who develop opioid use disorder, the prognosis improves significantly with early access to multidisciplinary support, including opioid agonist therapy (such as methadone or buprenorphine) and cognitive behavioural therapy.

\section*{Prevention and Self-Care}
\begin{itemize}
    \item Take opioid medicines strictly as prescribed by a medical practitioner, and never adjust the dose or frequency without consultation.
    \item Avoid combining opioids with alcohol, benzodiazepines, or other sedating medicines, as this dramatically increases the risk of fatal overdose.
    \item Store opioid medications securely in a locked cabinet, away from children and pets, and dispose of unused medicines at a local pharmacy.
    \item Maintain hydration and a high-fibre diet, and use prescribed laxatives to prevent opioid-induced constipation.
    \item Participate actively in a comprehensive pain management plan, incorporating physical therapy and non-pharmacological strategies.
\end{itemize}

\section*{Sources and Further Reading}
The following reputable organisations provide further information and clinical guidelines on the pharmacology and clinical use of opioid medicines:
\begin{itemize}
    \item Therapeutic Guidelines (Australia). \textit{Analgesic guidelines and opioid prescribing principles.} https://www.tg.org.au/
    \item NPS MedicineWise. \textit{Resources on opioid medicines, dosage, and safety.} https://www.nps.org.au/
    \item Faculty of Pain Medicine, ANZCA. \textit{Clinical statements and guidelines on opioid therapy.} https://www.fpm.anzca.edu.au/
    \item Mayo Clinic. \textit{Patient education on how opioids work and their risks.} https://www.mayoclinic.org/
    \item World Health Organization. \textit{Guidelines for the pharmacological and nutritional management of pain.} https://www.who.int/
\end{itemize}
